\begin{proposition}
\label{proposition-ML-tensor}
Let $M$ be an $R$-module.  The following are equivalent:
\begin{enumerate}
\item $M$ is Mittag-Leffler.
\item For every family $(Q_{\alpha})_{\alpha \in A}$ of $R$-modules, the
canonical map $M \otimes_R \left( \prod_{\alpha} Q_{\alpha} \right)
\to \prod_{\alpha} (M \otimes_R Q_{\alpha})$ is injective.
\end{enumerate}
\end{proposition}

\begin{proof}
First we prove (1) implies (2).  Suppose $M$ is Mittag-Leffler and let $x$ be
in the kernel of $M \otimes_R (\prod_{\alpha} Q_{\alpha}) \to
\prod_{\alpha} (M \otimes_R Q_{\alpha})$.  Write $M$ as a colimit $M =
\colim_{i \in I} M_i$ of a directed system of finitely presented modules
$M_i$.
 Then $M \otimes_R (\prod_{\alpha} Q_{\alpha})$ is the colimit of $M_i
\otimes_R (\prod_{\alpha} Q_{\alpha})$.  So $x$ is the image of an element
$x_i \in M_i \otimes_R (\prod_{\alpha} Q_{\alpha})$.  We must show that $x_i$
maps to $0$ in $M_j \otimes_R (\prod_{\alpha} Q_{\alpha})$ for some $j \geq
i$.  Since $M$ is Mittag-Leffler, we may choose $j \geq i$ such that $M_i
\to M_j$ and $M_i \to M$ dominate each other.  Then consider
the commutative diagram
$$
\xymatrix{
M \otimes_R (\prod_{\alpha} Q_{\alpha}) \ar[r] & \prod_{\alpha} (M
\otimes_R Q_{\alpha}) \\
M_i \otimes_R (\prod_{\alpha} Q_{\alpha}) \ar[r]^{\cong} \ar[d] \ar[u] &
\prod_{\alpha} (M_i \otimes_R Q_{\alpha}) \ar[d] \ar[u] \\
M_j \otimes_R (\prod_{\alpha} Q_{\alpha}) \ar[r]^{\cong} & \prod_{\alpha}
(M_j \otimes_R Q_{\alpha})
}
$$
whose bottom two horizontal maps are isomorphisms, according to Proposition
\ref{proposition-fp-tensor}.  Since $x_i$ maps to $0$ in $\prod_{\alpha} (M
\otimes_R Q_{\alpha})$, its image in $\prod_{\alpha} (M_i \otimes_R
Q_{\alpha})$ is in the kernel of the map $\prod_{\alpha} (M_i \otimes_R
Q_{\alpha}) \to \prod_{\alpha} (M \otimes_R Q_{\alpha})$.  But this
kernel equals the kernel of $\prod_{\alpha} (M_i \otimes_R Q_{\alpha})
\to \prod_{\alpha} (M_j \otimes_R Q_{\alpha})$ according to the
choice of $j$.  Thus $x_i$ maps to $0$ in $\prod_{\alpha} (M_j \otimes_R
Q_{\alpha})$ and hence to $0$ in $M_j \otimes_R (\prod_{\alpha} Q_{\alpha})$.

\medskip\noindent
Now suppose (2) holds. We prove $M$ satisfies formulation (1) of being
Mittag-Leffler from Proposition \ref{proposition-ML-characterization}.  Let $f:
P \to M$ be a map from a finitely presented module $P$ to $M$.  Choose
a set $B$ of representatives of the isomorphism classes of finitely presented
$R$-modules. Let $A$ be the set of pairs $(Q, x)$ where $Q \in B$ and $x \in
\Ker(P \otimes Q \to M \otimes Q)$.  For $\alpha = (Q, x) \in A$, we
write $Q_{\alpha}$ for $Q$ and $x_{\alpha}$ for $x$.  Consider the commutative
diagram
$$
\xymatrix{
M \otimes_R (\prod_{\alpha} Q_{\alpha}) \ar[r] &
\prod_{\alpha} (M \otimes_R Q_{\alpha}) \\
P \otimes_R (\prod_{\alpha} Q_{\alpha}) \ar[r]^{\cong} \ar[u] &
\prod_{\alpha} (P \otimes_R Q_{\alpha}) \ar[u] .
}
$$
The top arrow is an injection by assumption, and the bottom arrow is an
isomorphism by Proposition \ref{proposition-fp-tensor}.  Let $x \in P
\otimes_R (\prod_{\alpha} Q_{\alpha})$ be the element corresponding to
$(x_{\alpha}) \in \prod_{\alpha} (P \otimes_R Q_{\alpha})$ under this
isomorphism.  Then $x \in \Ker( P \otimes_R (\prod_{\alpha} Q_{\alpha})
\to M \otimes_R (\prod_{\alpha} Q_{\alpha}))$ since the top arrow in
the diagram is injective.  By Lemma \ref{lemma-kernel-tensored-fp}, we get a
finitely presented module $P'$ and a map $f': P \to P'$ such that $f: P
\to M$ factors through $f'$ and $x \in \Ker(P \otimes_R
(\prod_{\alpha} Q_{\alpha}) \to P' \otimes_R (\prod_{\alpha}
Q_{\alpha}))$.  We have a commutative diagram
$$
\xymatrix{
P' \otimes_R (\prod_{\alpha} Q_{\alpha}) \ar[r]^{\cong} &
\prod_{\alpha} (P' \otimes_R Q_{\alpha}) \\
P \otimes_R (\prod_{\alpha} Q_{\alpha}) \ar[r]^{\cong} \ar[u] &
\prod_{\alpha} (P \otimes_R Q_{\alpha}) \ar[u] .
}
$$
where both the top and bottom arrows are isomorphisms by Proposition
\ref{proposition-fp-tensor}.  Thus since $x$ is in the kernel of the left
vertical map, $(x_{\alpha})$ is in the kernel of the right vertical map.  This
means $x_{\alpha} \in \Ker(P \otimes_R Q_{\alpha} \to P' \otimes_R
Q_{\alpha})$ for every $\alpha \in A$.  By the definition of $A$ this means
$\Ker(P \otimes_R Q \to P' \otimes_R Q) \supset \Ker(P \otimes_R
Q \to M \otimes_R Q)$ for all finitely presented $Q$ and, since $f: P
\to M$ factors through $f': P \to P'$, actually equality holds.
 By Lemma \ref{lemma-domination-fp}, $f$ and $f'$ dominate each other.
\end{proof}